\chapter{DESCRIÇÃO DO LOCAL DE ESTÁGIO}
\label{cap:local}

A Universidade Federal de Minas Gerais está localizada na Avenida
Presidente Antônio Carlos, n° 6627, no bairro Pampulha, na cidade de
Belo Horizonte, estado de Minas Gerais. O Hospital Veterinário é um
órgão complementar à Escola de Veterinária da UFMG que desenvolve
atividades de ensino, pesquisa e extensão, composto pelos setores de
Clínica Médica, Clínica Cirúrgica, Patologia, Reprodução e Divisão de
Enfermagem. Nesse contexto, realizam-se atendimentos a diferentes
espécies de animais domésticos, bem como a algumas espécies silvestres
\cite{UFMG2026}.

A Clínica Médica de Equinos da UFMG é o setor responsável por receber
atendimentos dos mais variados na medicina equina, atuando em urgências,
emergências e ainda atendimentos eletivos. A referida encontra-se
instalada no Galpão D do Hospital Veterinário, com funcionamento regular
das 8h às 18h, sendo realizados plantões noturnos durante a semana e aos
finais de semana, conforme a necessidade do serviço, a fim de assegurar
a continuidade da assistência aos animais atendidos. A equipe fixa do
setor é composta por quatro médicos veterinários residentes, três
professoras coordenadoras, uma enfermeira, médicos veterinários
plantonistas e tratadores, além de estagiários vinculados ao estágio
supervisionado obrigatório e à vivência acadêmica.

\section{Estrutura física}
\label{sec:estrutura}

As instalações do setor de Clínica Médica de Equinos da UFMG são bem
estruturadas e atendem às demandas da rotina clínica hospitalar. Na área
interna, o setor dispõe de uma farmácia, três troncos de contenção, uma
bancada principal equipada com itens de uso rotineiro, doze baias, uma
sala destinada aos residentes e aos professores coordenadores, um espaço
voltado aos estagiários que também desempenha função de controle de
itens, uma sala de equipamentos, uma copa e um quarto de rações. Na área
externa, encontram-se um redondel, um desembarcador, um galpão destinado
ao armazenamento de feno e um galpão denominado ``carroceiros'', o qual
dispõe de três baias e dois troncos de contenção, utilizados conforme a
demanda do setor. Além disso, a estrutura conta com dezesseis piquetes,
sendo quinze equipados com cochos de alimentação e bebedouros.

\subsection{Farmácia}
\label{subsec:farmacia}

O setor conta com duas farmácias para suporte à rotina clínica. A
primeira é de uso compartilhado entre os diferentes setores do Hospital
Veterinário, encontrando-se instalada no setor de pequenos animais. A
segunda, inaugurada no galpão durante as primeiras semanas do período de
estágio supervisionado, atende exclusivamente aos setores de Clínica
Cirúrgica de Grandes Animais e Clínica Médica de Equinos. Ambas unidades
estão abastecidas com os medicamentos e insumos necessários ao
atendimento clínico, sendo os itens solicitados conforme a demanda pelos
médicos veterinários residentes e docentes responsáveis.

\begin{figure}[htb]
  \centering
  \caption{Farmácia de uso comum do Hospital Veterinário, EV-UFMG.}
  \label{fig:farmacia1}
  \includegraphics[width=0.5\textwidth]{media/image6.jpg}
  \fonte{Da autora (2026).}
\end{figure}

\begin{figure}[htb]
  \centering
  \caption{Farmácia destinada aos setores de Clínica Cirúrgica de Grandes Animais e Clínica Médica de Equinos, EV-UFMG.}
  \label{fig:farmacia2}
  \includegraphics[width=0.5\textwidth]{media/image8.jpg}
  \fonte{Da autora (2026).}
\end{figure}

\subsection{Sala dos Residentes}
\label{subsec:sala-residentes}

A sala dos residentes é o espaço de uso predominante por médicos
veterinários residentes e docentes coordenadores, sendo destinada
principalmente às atividades administrativas. Nesse ambiente,
realizam-se reuniões e discussões de casos clínicos, além de constituir
um espaço de organização individual e coletiva da equipe. Ademais,
abriga o computador principal do setor, por meio do qual se acessa o
sistema próprio do Hospital Veterinário.

\begin{figure}[htb]
  \centering
  \caption{Sala dos residentes do setor de Clínica Médica de Equinos da EV-UFMG.}
  \label{fig:sala-residentes}
  % [INSERIR IMAGEM]
  \fonte{Da autora (2026).}
\end{figure}

\subsection{Sala dos Estagiários}
\label{subsec:sala-estagiarios}

A sala dos estagiários, como o próprio nome sugere, é destinada
predominantemente aos estagiários, sejam eles vinculados ao estágio
supervisionado ou de vivência, constituindo um espaço designado ao
armazenamento de pertences pessoais e à utilização das mesas para
estudo. Além disso, constitui também um ambiente de apoio à rotina
clínica, no qual os materiais encontram-se organizados em quatro
armários e uma estante, sendo cada um destinado a uma função específica,
abrangendo desde itens de uso geral até aqueles voltados a emergências.

\begin{figure}[htb]
  \centering
  \caption{Sala dos estagiários do setor de Clínica Médica de Equinos, EV-UFMG.}
  \label{fig:sala-estagiarios}
  \includegraphics[width=0.55\textwidth]{media/image15.jpg}
  \fonte{Da autora (2026).}
\end{figure}

\subsection{Sala de Equipamentos}
\label{subsec:sala-equipamentos}

A sala de equipamentos destina-se ao armazenamento de equipamentos de
uso específico, como equipamentos de ultrassonografia e endoscopia, além
de materiais destinados ao atendimento a campo.

\begin{figure}[htb]
  \centering
  \caption{Sala de equipamentos do setor de Clínica Médica de Equinos, EV-UFMG.}
  \label{fig:sala-equipamentos}
  \includegraphics[width=0.5\textwidth]{media/image18.jpg}
  \fonte{Da autora (2026).}
\end{figure}

\subsection{Baias}
\label{subsec:baias}

A Clínica Médica de Equinos dispõe de 12 baias na área interna,
distribuídas em dois corredores, que medem cada uma aproximadamente 4
metros de comprimento por 4 metros de largura, destinadas à recepção,
acomodação e cuidados dos pacientes. Cada baia é constituída por portas
e janelas de madeira, além de dois cochos destinados à alimentação e um
para fornecimento de água. Nas portas, encontram-se afixadas fichas de
identificação que contém informações essenciais para a individualização
dos pacientes, bem como orientações de manejo, como quantidade de
concentrado e volumoso fornecidos, altura de cama e necessidade de
jejum. Dentre as doze baias, uma destina-se especificamente ao manejo de
animais suspeitos de doenças infectocontagiosas, localizada ao final do
corredor, com o objetivo de favorecer o isolamento e a biossegurança do
setor.

\begin{figure}[htb]
  \centering
  \caption{Corredor com baias do setor de Clínica Médica de Equinos, EV-UFMG.}
  \label{fig:corredor-baias}
  \includegraphics[width=0.55\textwidth]{media/image26.jpg}
  \fonte{Da autora (2026).}
\end{figure}

\begin{figure}[htb]
  \centering
  \caption{Vista interna de baia do setor de Clínica Médica de Equinos, EV-UFMG.}
  \label{fig:baia-interna}
  \includegraphics[width=0.55\textwidth]{media/image17.jpg}
  \fonte{Da autora (2026).}
\end{figure}

\begin{figure}[htb]
  \centering
  \caption{Ficha de identificação do paciente afixada na porta da baia, EV-UFMG.}
  \label{fig:ficha-baia}
  \includegraphics[width=0.55\textwidth]{media/image10.jpg}
  \fonte{Da autora (2026).}
\end{figure}

\begin{figure}[htb]
  \centering
  \caption{Baia destinada ao isolamento de pacientes suspeitos de doenças infectocontagiosas, setor de Clínica Médica de Equinos, EV-UFMG.}
  \label{fig:baia-isolamento}
  \includegraphics[width=0.5\textwidth]{media/image11.jpg}
  \fonte{Da autora (2026).}
\end{figure}

\subsection{Troncos de Contenção}
\label{subsec:troncos}

A Clínica dispõe de três troncos de contenção, sendo dois localizados na
área principal, destinados à recepção dos animais no momento da
admissão, onde são realizados os exames clínicos iniciais. Esses troncos
também são utilizados para a execução de procedimentos, como
fluidoterapia, bem como para a realização de exames complementares, como
a ultrassonografia, entre outros. O terceiro tronco encontra-se em área
adjacente, sendo utilizado predominantemente em atividades didáticas;
contudo, pode ser empregado na rotina clínica conforme a demanda do
setor.

\begin{figure}[htb]
  \centering
  \caption{Troncos de contenção da área principal do setor de Clínica Médica de Equinos, EV-UFMG.}
  \label{fig:tronco-principal}
  \includegraphics[width=0.5\textwidth]{media/image20.jpg}
  \fonte{Da autora (2026).}
\end{figure}

\begin{figure}[htb]
  \centering
  \caption{Tronco de contenção da área adjacente do setor de Clínica Médica de Equinos, EV-UFMG.}
  \label{fig:tronco-adjacente}
  \includegraphics[width=0.5\textwidth]{media/image13.jpg}
  \fonte{Da autora (2026).}
\end{figure}

Adjacente aos troncos da área principal, encontra-se uma bancada
equipada com pia, destinada à higienização das mãos e dos materiais
utilizados nos atendimentos. Sobre a bancada são mantidos materiais de
uso rotineiro, incluindo produtos para antissepsia, luvas, pomadas,
limpadores de casco, escovas e demais insumos necessários ao suporte dos
procedimentos realizados na área.

\begin{figure}[htb]
  \centering
  \caption{Bancada com pia e materiais de uso rotineiro, EV-UFMG.}
  \label{fig:bancada}
  \includegraphics[width=0.5\textwidth]{media/image14.jpg}
  \fonte{Da autora (2026).}
\end{figure}

\subsection{Quarto de ração}
\label{subsec:quarto-racao}

O setor conta com um quarto de rações de uso comum à Clínica Médica de
Equinos e à Clínica Cirúrgica de Grandes Animais, onde são armazenados,
em bombonas devidamente vedadas, os concentrados utilizados na rotina
alimentar dos pacientes, bem como feno peletizado e sal mineral
específico para equinos.

\begin{figure}[htb]
  \centering
  \caption{Quarto de rações de uso comum à Clínica Médica de Equinos e à Clínica Cirúrgica de Grandes Animais, EV-UFMG.}
  \label{fig:quarto-racao}
  \includegraphics[width=0.55\textwidth]{media/image19.jpg}
  \fonte{Da autora (2026).}
\end{figure}

\subsection{Piquetes}
\label{subsec:piquetes}

A área externa é composta por dezesseis piquetes de dimensões variadas.
Um deles está localizado ao lado do setor e é delimitado por correntes,
sendo utilizado para a soltura momentânea dos animais para pastejo. Os
demais, situados a maior distância, são cercados com arame liso,
garantindo a contenção adequada dos equinos. Estes dispõem de bebedouros
e cochos para alimentação, além de cobertura que protege tanto os
animais quanto o alimento das variações climáticas. Em relação à
capacidade de lotação, parte dos piquetes é destinada à acomodação
individual, ao passo que outros comportam até dois equinos adultos.

\begin{figure}[htb]
  \centering
  \caption{Piquetes da área externa do setor de Clínica Médica de Equinos, EV-UFMG.}
  \label{fig:piquetes}
  \includegraphics[width=0.9\textwidth]{media/image28.jpg}
  \fonte{Da autora (2026).}
\end{figure}

\subsection{Desembarcador}
\label{subsec:desembarcador}

O desembarcador constitui o ponto de entrada e saída dos animais no
setor, sendo utilizado tanto para a recepção dos pacientes no momento da
admissão quanto para o embarque daqueles que receberam alta hospitalar.
Trata-se de uma estrutura fundamental para a segurança no manejo dos
equinos durante o transporte, permitindo que o desembarque e o embarque
sejam realizados de forma adequada e com menor risco de injúrias tanto
aos animais quanto aos profissionais envolvidos.

\begin{figure}[htb]
  \centering
  \caption{Desembarcador do Hospital Veterinário, EV-UFMG.}
  \label{fig:desembarcador}
  \includegraphics[width=0.5\textwidth]{media/image23.jpg}
  \fonte{Da autora (2026).}
\end{figure}

\subsection{Redondel}
\label{subsec:redondel}

O setor conta com dois redondeis, um fechado e um aberto, estruturas
circulares destinadas à movimentação e avaliação dos equinos em ambiente
controlado. Sua principal finalidade é permitir a observação do sistema
locomotor dos animais e também constitui um espaço para o
desenvolvimento de protocolos de fisioterapia e reabilitação dos
pacientes. Além disso, o redondel fechado é utilizado, quando
necessário, como espaço para a soltura dos animais, proporcionando uma
alternativa de movimentação livre fora da baia.

\begin{figure}[htb]
  \centering
  \caption{Redondel fechado do Hospital Veterinário, EV-UFMG.}
  \label{fig:redondel}
  \includegraphics[width=0.65\textwidth]{media/image27.jpg}
  \fonte{Da autora (2026).}
\end{figure}

\subsection{Galpão de Feno}
\label{subsec:galpao-feno}

O galpão de feno destina-se ao armazenamento do volumoso utilizado na
alimentação dos pacientes internados. A estrutura coberta e abrigada
protege o alimento das intempéries, preservando sua qualidade
nutricional.

\begin{figure}[htb]
  \centering
  \caption{Galpão de feno da área externa do Hospital Veterinário, EV-UFMG.}
  \label{fig:galpao-feno}
  \includegraphics[width=0.5\textwidth]{media/image29.jpg}
  \fonte{Da autora (2026).}
\end{figure}

\subsection{Galpão de Carroceiros}
\label{subsec:galpao-carroceiros}

O galpão de carroceiros é destinado predominantemente ao desenvolvimento
de atividades didáticas e de pesquisa, sendo utilizado em aulas práticas
e experimentos de iniciação científica. A estrutura conta com três baias
e dois troncos de contenção, podendo ser acionada como suporte à rotina
hospitalar em momentos de maior demanda. Dispõe ainda de uma sala
destinada à estocagem de ração, destinada ao armazenamento de insumos
excedentes que não são estocados na farmácia principal do setor, como
sacos fechados de ração, e sal, garantindo a organização e conservação
adequada dos insumos utilizados no manejo alimentar dos animais.

\begin{figure}[htb]
  \centering
  \caption{Galpão de carroceiros do setor de Clínica Médica de Equinos, EV-UFMG. Troncos de contenção (A). Baias (B).}
  \label{fig:galpao-carroceiros}
  \includegraphics[width=0.45\textwidth]{media/image21.jpg}
  \includegraphics[width=0.45\textwidth]{media/image25.jpg}
  \fonte{Da autora (2026).}
\end{figure}
